% EXPECTED-LEAKS: printanswers, question, choice, CorrectChoice, solution, fillwithlines, bonuspoints, pointsofquestion, qformat
\documentclass[12pt,answers]{exam}
\usepackage{amsmath}

\printanswers

\begin{document}

\begin{center}
{\Large Midterm Exam --- CS 101} \\
\vspace{0.2cm}
Name: \rule{5cm}{0.2pt} \hfill Date: \rule{3cm}{0.2pt}
\end{center}

\begin{questions}

\question[10] What is the time complexity of binary search on a sorted array of size $n$?

\begin{choices}
  \choice $O(n)$
  \choice $O(n^2)$
  \CorrectChoice $O(\log n)$
  \choice $O(1)$
\end{choices}

\begin{solution}
Binary search halves the search space at each step, giving logarithmic complexity.
\end{solution}

\question[15] Prove that the sum $1 + 2 + \cdots + n = \frac{n(n+1)}{2}$ for all positive integers $n$.

\begin{solution}
By induction on $n$. Base case: $n = 1$ gives $1 = 1(2)/2$. Inductive step: assume the formula holds for $k$; then $1 + 2 + \cdots + k + (k+1) = \frac{k(k+1)}{2} + (k+1) = \frac{(k+1)(k+2)}{2}$.
\end{solution}

\question[25] Implement quicksort in pseudocode. Analyse the worst-case and expected runtime.

\fillwithlines{3in}

\end{questions}

\end{document}
